%% Credal Concept Formation in Dependent Type Theory
%%
%% Render with:
%%   lualatex -shell-escape credal-concept-formation.tex \
%%     && bibtex credal-concept-formation \
%%     && lualatex -shell-escape credal-concept-formation.tex \
%%     && lualatex -shell-escape credal-concept-formation.tex

\documentclass[]{article}
\usepackage{url}
\usepackage[hyperindex,breaklinks]{hyperref}
\IfFileExists{breakurl.sty}{\usepackage{breakurl}}{}
\usepackage{listings}
\lstset{basicstyle=\ttfamily\footnotesize,breaklines=true,frame=single,
        xleftmargin=1em,xrightmargin=0.5em}
\usepackage{float}
\restylefloat{table}
\usepackage{longtable}
\usepackage{booktabs}
\usepackage{graphicx}
\usepackage[font=small,labelfont=bf]{caption}
\usepackage[skip=0pt]{subcaption}
\usepackage{amsfonts}
\usepackage{amsmath}
\usepackage{amsthm}
\usepackage{unicode-math}
\IfFileExists{newunicodechar.sty}{
  \usepackage{newunicodechar}
  \newunicodechar{∘}{\ensuremath{\circ}}
}{}
\IfFontExistsTF{Stix Two Math}{\setmathfont{Stix Two Math}}{\setmathfont{STIX Math}}

\makeatletter
\providecommand{\headerps@out}{}
\makeatother

\newtheorem{theorem}{Theorem}
\newtheorem{lemma}[theorem]{Lemma}
\newtheorem{proposition}[theorem]{Proposition}
\newtheorem{corollary}[theorem]{Corollary}
\newtheorem{definition}[theorem]{Definition}
\newtheorem{remark}[theorem]{Remark}
\newtheorem{observation}[theorem]{Observation}

\newcommand{\code}[1]{\texttt{#1}}
\newcommand{\codepath}[1]{\path{#1}}
\newcommand{\lid}[1]{{\ttfamily\nolinkurl{#1}}}
\newcommand{\lpath}[1]{\texttt{\nolinkurl{#1}}}
\newcommand{\ty}[1]{\ensuremath{\mathtt{#1}}}
\newcommand{\STV}[2]{\langle #1,\, #2 \rangle}
\newcommand{\strength}{\mathrm{strength}}
\newcommand{\confidence}{\mathrm{confidence}}
\newcommand{\Eq}{\;\Leftrightarrow\;}
\newcommand{\LowerFC}{\mathcal{L}}
\newcommand{\UpperFC}{\mathcal{U}}
\newcommand{\PreciseEnv}{\mathcal{P}}

\emergencystretch=3em

\title{Credal Concept Formation in Dependent Type Theory:\\
       FCA Recovery, Walley Conjugacy, and a Mizar Benchmark}
\author{Mettapedia Working Notes}
\date{\today}

\begin{document}
\maketitle

\begin{abstract}
We give a type-theoretic formalization of credal concept formation in Lean~4 that admits
classical Formal Concept Analysis (FCA) as an exact, order-isomorphic special case.
The contributions are: (i) subtype-based lower and upper formed-concept carriers,
\lid{LowerFormedConcept} and \lid{UpperFormedConcept}, parameterized over a family of
admissible evidence gates; (ii) an order-isomorphism recovery theorem
\lid{formedConceptOrderIsoCrispConcept} establishing that our exact lane \emph{is} the
classical FCA concept lattice; (iii) singleton-collapse equivalences that prove
conservative extension of the crisp case; (iv) integration with Walley's natural
extension via a four-scale conjugacy cluster on a projective credal spec, lifting the
credal coordinates (lower, upper, midpoint, width-complement) into the
imprecise-probability machinery the PLN $\STV{s}{c}$ truth value compresses; (v) a
first-stone empirical benchmark over an extracted slice of the Mizar Mathematical
Library article \lid{conlat\_1.miz}, including a real credal witness on the
\lid{ObjectDerivation} base concept. All theorems are quoted verbatim from a public
Lean~4 source with no \lid{sorry}, \lid{admit}, or new \lid{axiom}. We are explicit
about what is theorem-witness and what is yet-to-run experiment; \S 9 enumerates the
empirical work required to turn the benchmark into a published discovery.
\end{abstract}

\section{Introduction}\label{sec:intro}

Classical Formal Concept Analysis (FCA), introduced by Wille~\cite{wille1982}, gives a
crisp, order-theoretic account of how concepts arise from an incidence relation
$M : \mathrm{Obj} \times \mathrm{Attr} \to \{0,1\}$. The standard objects are
\emph{formal concepts}: pairs $(A, B)$ with $A = B^\downarrow$ and $B = A^\uparrow$
under the derivation operators. The standard structure is a complete lattice with
intent/extent duality. After forty years and several thousand papers, this is the
canonical mathematical framework for the question ``what concepts arise from this
context?'' For practical concept formation over noisy or evidence-graded data,
however, the binary commitment ${M : \mathrm{Obj} \times \mathrm{Attr} \to \{0,1\}}$ is
too rigid.

Several lines of work generalize FCA to handle uncertainty. Belohlavek's fuzzy
FCA~\cite{belohlavek1999} replaces the crisp incidence with degrees in $[0,1]$.
Djouadi and Prade introduce interval-valued and possibilistic
contexts~\cite{djouadi2009}. Triadic FCA~\cite{wille1995} adds a context dimension.
Each gives one principled answer to ``what if the relation $M$ is uncertain?''.

This paper takes a different route. We treat the uncertainty as living \emph{not in
the relation $M$ itself}, but in which \emph{admissible evidence gate} converts
evidence-valued cells into crisp incidence assertions. A gate is a thresholding
discipline; admissible gates are those the modeller is willing to consider; concept
formation under \emph{every} admissible gate gives a \emph{lower} carrier
(``robustly formed concepts'') while concept formation under \emph{some} admissible
gate gives an \emph{upper} carrier (``permissively formed concepts''). This matches
the Walley imprecise-probability framework~\cite{walley1991}: the gate family
indexes a family of precise previsions; the lower/upper concepts are the credal
envelopes of an underlying gamble. The exact (singleton-gate) case recovers
classical FCA.

\paragraph{Three concrete contributions.} The paper formalizes this picture in
Lean~4 and ties it to existing PLN infrastructure as follows.

\begin{enumerate}
\item \textbf{Type-theoretic credal carriers (\S\ref{sec:carriers}).} We define
\lid{LowerFormedConcept} and \lid{UpperFormedConcept} as proper subtypes of
\lid{DualConcept}, not as wrappers around \lid{FormedConcept}. The subtype design
avoids a class of degenerate-instance traps; the upper-side witness extraction
uses \lid{Classical.choose} honestly.

\item \textbf{Order-isomorphic FCA recovery (\S\ref{sec:fca-recovery}).} We prove
\lid{formedConceptEquivCrispConcept} (the carrier equivalence) and
\lid{formedConceptOrderIsoCrispConcept} (the order isomorphism) against mathlib's
\lid{CrispConcept}. The recovery is not a structural resemblance: the exact lane
\emph{is} the classical FCA concept lattice as an order-theoretic object.

\item \textbf{Walley natural-extension bridge (\S\ref{sec:walley}).} Four
parameterized conjugacy theorems at the cylinder, finite-cylinder, global, and
finite-global scales of a projective local credal spec connect upper-envelope and
natural-extension semantics. The credal coordinates \lid{credalEnvelopeWidth},
\lid{credalEnvelopeMidpoint}, \lid{credalEnvelopeWidthComplement} are the natural
$\STV{s}{c}$-compression coordinates the PLN truth value rides on.
\end{enumerate}

\paragraph{Empirical anchor.} \S\ref{sec:mizar} reports a first-stone benchmark: a
Python extractor walks a \lid{.miz} article and emits an attribute-counted
observation context; \S\ref{sec:mizar} exhibits a real credal witness on the
\lid{ObjectDerivation} base concept of \lid{conlat\_1.miz} where the concept lives in
the upper carrier but not the lower, yielding $\mathrm{width} = 1,\
\mathrm{midpoint} = \tfrac{1}{2}$. We are explicit (\S\ref{sec:mizar-honest}) that this
is theorem-witness over hand-chosen gates and one article. \S\ref{sec:roadmap}
discusses what real experiments would look like, including the most interesting one:
using the upper/lower gap on a multi-article Mizar corpus to detect
\emph{formally-missing concepts} via Galois-dual reasoning.

\section{Background}\label{sec:background}

\subsection{Classical FCA}\label{sec:bg-fca}

A formal context is a triple $(\mathrm{Obj}, \mathrm{Attr}, I)$ with
$I \subseteq \mathrm{Obj} \times \mathrm{Attr}$. The derivation operators
$\cdot^\uparrow : 2^{\mathrm{Obj}} \to 2^{\mathrm{Attr}}$ and
$\cdot^\downarrow$ form a Galois connection. A formal concept is a pair
$(A, B)$ with $A^\uparrow = B$ and $B^\downarrow = A$. The set of formal concepts
ordered by $(A, B) \leq (A', B') \iff A \subseteq A'$ is the
\emph{concept lattice} of the context. The main theorem of
FCA~\cite{wille1982,gantnerwille1999} is that every finite lattice arises this way.

Algorithmic FCA tooling (NextClosure, ConExp, FCAbedrock, Lattice Miner) computes
concept lattices in time proportional to lattice size, not $2^{|\mathrm{Obj}|}$.
For practical contexts the lattice stays small. The crisp commitment, however, is
brittle: a single attribute-cell flip can move concepts in and out of the lattice
discontinuously.

\subsection{Imprecise probability and Walley's natural extension}\label{sec:bg-walley}

Walley~\cite{walley1991} replaces a single probability measure with a \emph{credal
set} of admissible precise previsions. A lower prevision $\underline{P}$ is
\emph{coherent} when it avoids sure loss and satisfies superadditivity and positive
homogeneity. The \emph{natural extension} of a coherent lower prevision is the
greatest coherent lower prevision dominated by it; its conjugate $\overline{P}(X)
= -\underline{P}(-X)$ is the corresponding upper prevision. The
\emph{lower envelope} of a credal set $\mathcal{C}$ on a gamble $X$ is
$\inf_{P \in \mathcal{C}} P(X)$; the \emph{upper envelope} is $\sup$. Walley's
representation theorem links the two notions: every coherent lower prevision is the
lower envelope of its set of dominating precise previsions, and conversely.

Belohlavek's fuzzy FCA~\cite{belohlavek1999} and Djouadi--Prade's interval-valued
FCA~\cite{djouadi2009} are predecessors for treating FCA-style closures under
uncertainty. They do not bridge into Walley's natural-extension calculus, and do
not formalize a conservative reduction to crisp FCA in dependent type theory.

\subsection{PLN truth values}\label{sec:bg-pln}

Probabilistic Logic Networks~\cite{goertzel2008} represent the strength and
confidence of an inheritance relation as a two-number tuple $\STV{s}{c}$, where
$s \in [0,1]$ is interpreted as a strength-like point estimate and
$c \in [0,1]$ as a confidence-like inverse-imprecision. A companion characterization
paper~\cite{pln-credal-degrees-of-freedom} identifies $\STV{s}{c}$ as a two-number
compression of an underlying credal set, with three layered slots where the
compression admits genuine degrees of freedom. The present paper provides the
\emph{concept-formation} side of that picture: how credal carriers, FCA recovery,
and Walley conjugacy together justify a midpoint/width-complement reading of
$\STV{s}{c}$.

\section{The credal carriers}\label{sec:carriers}

The exact (precise) concept-formation surface in our framework is given by the
\lid{ObservationSurface} machinery and its
\lid{finiteConceptFamily} operator; see~\cite{wmpln_concept_formation_guidance} for
the public-API discussion. Codex's \lid{ConceptOntology/Formation.lean} provides the
operator and a correctness theorem
\(\,A \in \operatorname{finiteConceptFamily}(S,G,\sigma)\,\) iff \(A\) is closed
under \(\operatorname{crispRelation}(G,S.\operatorname{aggregate}(\sigma))\).
The credal extension lives one level up.

\paragraph{Carriers as subtypes.} The credal carriers are not wrappers; they are
subtypes whose proof field carries genuine content.

\begin{definition}[\lid{LowerFormedConcept} and \lid{UpperFormedConcept},
\lpath{ConceptOntology/CredalFormation.lean:143,148}]\label{def:carriers}
\begin{lstlisting}
abbrev LowerFormedConcept
    (Γ : Gate → EvidenceGate Q) (M : Obj → Attr → Q) :=
  { A : DualConcept Obj Attr // A ∈ lowerConceptFamily Γ M }

abbrev UpperFormedConcept
    (Γ : Gate → EvidenceGate Q) (M : Obj → Attr → Q) :=
  { A : DualConcept Obj Attr // A ∈ upperConceptFamily Γ M }
\end{lstlisting}
\end{definition}

\noindent The membership predicate \lid{lowerConceptFamily} is ``$A$ is a formed
concept under \emph{every} gate $\Gamma(g)$'' and \lid{upperConceptFamily} is ``$A$
is a formed concept under \emph{some} gate.'' The carriers' proof field \(A.2\)
certifies \(A \in \operatorname{lowerConceptFamily}(\Gamma,M)\): an honest
dependent proof obligation, not a definitional unfolding.

\paragraph{Why subtype, not wrapper.} The carrier design avoids a degenerate
pattern that would render the credal layer vacuous: defining a ``credal concept''
as a \emph{renamed} \lid{FormedConcept} and proving the recovery theorem by
\lid{rfl}. The present subtype design forces the recovery theorem to do real work
because the carrier types are genuinely distinct, even though their meaning sets are
related.

\paragraph{Working with the proof field.} The single-gate projection
\lid{lowerFormedConceptAt} uses the universal-gate proof concretely:

\begin{lstlisting}
def lowerFormedConceptAt
    (Γ : Gate → EvidenceGate Q) (M : Obj → Attr → Q)
    (g : Gate) :
    LowerFormedConcept Γ M → AbstractInheritance.FormedConcept (Γ g) M :=
  fun A => ⟨A.1, A.2 g⟩
\end{lstlisting}

\noindent (\lpath{ConceptOntology/CredalFormation.lean:198}.) The upper-side
witness extraction uses \lid{Classical.choose}:

\begin{lstlisting}
noncomputable def upperFormedConceptWitness
    (Γ : Gate → EvidenceGate Q) (M : Obj → Attr → Q) :
    UpperFormedConcept Γ M →
      Σ g : Gate, AbstractInheritance.FormedConcept (Γ g) M := by
  classical
  intro A
  refine ⟨Classical.choose A.2, ?_⟩
  exact ⟨A.1, Classical.choose_spec A.2⟩
\end{lstlisting}

\noindent (line 205). Both projections preserve the underlying dual-concept value
unchanged; this is recorded as the \lid{lowerFormedConceptAt\_val} and
\lid{upperFormedConceptWitness\_val} simp lemmas.

\section{FCA recovery}\label{sec:fca-recovery}

Codex's \lpath{ConceptOntology/FCARecovery.lean} module makes the FCA recovery
explicit. The 74-line file proves two theorems and one corollary; we quote the
first verbatim and sketch the second.

\begin{theorem}[Carrier equivalence,
\lpath{ConceptOntology/FCARecovery.lean:27}]\label{thm:fca-carrier}
\begin{lstlisting}
noncomputable def formedConceptEquivCrispConcept
    (G : EvidenceGate Q) (M : Obj → Attr → Q) :
    AbstractInheritance.FormedConcept G M ≃ CrispConcept G M where
  toFun A :=
    DualConcept.toConcept A.1
      (AbstractInheritance.formedConcept_isClosed G M A)
  invFun c :=
    ⟨DualConcept.ofConcept c,
      (AbstractInheritance.mem_finiteConceptFamily_iff G M _).2
        (DualConcept.isClosed_ofConcept c)⟩
  left_inv A := by
    apply Subtype.ext
    exact DualConcept.ofConcept_toConcept A.1
      (AbstractInheritance.formedConcept_isClosed G M A)
  right_inv c := by
    exact DualConcept.toConcept_ofConcept c
\end{lstlisting}
\end{theorem}

\noindent The carrier equivalence promotes a formed-concept proof to a classical
formal-concept pair $(A, B)$ via the \lid{DualConcept.toConcept} construction, and
inverts it by \lid{DualConcept.ofConcept}. Both inverse laws are honest
sigma-extensionality lemmas, not \lid{rfl}.

\begin{theorem}[Order isomorphism,
\lpath{ConceptOntology/FCARecovery.lean:44}]\label{thm:fca-order}
\begin{lstlisting}
noncomputable def formedConceptOrderIsoCrispConcept
    (G : EvidenceGate Q) (M : Obj → Attr → Q) :
    AbstractInheritance.FormedConcept G M ≃o CrispConcept G M where
  toEquiv := formedConceptEquivCrispConcept G M
  map_rel_iff' := by
    intro A B
    simpa using
      (DualConcept.toConcept_le_iff A.1 B.1
        (AbstractInheritance.formedConcept_isClosed G M A)
        (AbstractInheritance.formedConcept_isClosed G M B))
\end{lstlisting}
\end{theorem}

\noindent This is the headline claim of the recovery layer: the exact formed-concept
order is the classical FCA concept-lattice order, as an order isomorphism. It is
not merely the case that the carriers are equipotent; the order structure transfers
exactly. A corollary at line 64,
\lid{formedConcept\_inherits\_iff\_crispConcept\_le}, packages the implication for
downstream consumers who want the inheritance-style direction.

\subsection{Singleton-collapse conservative extension}

The credal layer collapses to the exact layer in the singleton-gate case. We
formalize this via two type equivalences:

\begin{theorem}[Lower singleton equivalence,
\lpath{ConceptOntology/CredalFormation.lean:275}]\label{thm:lower-singleton}
\begin{lstlisting}
noncomputable def lowerFormedConceptSingletonEquiv
    (G : EvidenceGate Q) (M : Obj → Attr → Q) :
    LowerFormedConcept (Gate := PUnit) (fun _ => G) M ≃
      AbstractInheritance.FormedConcept G M
\end{lstlisting}
\end{theorem}

\noindent The companion \lid{upperFormedConceptSingletonEquiv} at line 292 is the
upper-side twin. Both proofs use the singleton iff lemmas
\lid{mem\_lowerConceptFamily\_singleton\_iff} and
\lid{mem\_upperConceptFamily\_singleton\_iff}, with \lid{cases A; rfl} discharging
the inverse laws. The constructions are honest equivalences between type-distinct
objects (the credal subtype \emph{with} \lid{Gate := PUnit} versus the bare
\lid{FormedConcept}).

\paragraph{What conservative extension means here.} When the gate family has a
single element, every \lid{LowerFormedConcept} is the same as a \lid{FormedConcept}
under that gate, and the same for \lid{UpperFormedConcept}. Composed with
Theorem~\ref{thm:fca-order}, this gives the chain
\[
  \text{Credal layer at singleton gate}
  \xrightarrow{\text{Thm.~\ref{thm:lower-singleton}}}
  \text{Exact \lid{FormedConcept}}
  \xrightarrow{\text{Thm.~\ref{thm:fca-order}}}
  \text{Classical FCA concept lattice}.
\]
Both arrows are order isomorphisms. The credal layer is a conservative refinement,
not a sideways alternative.

\section{Inheritance and transport}\label{sec:inheritance}

The credal carrier composes cleanly with the existing inheritance, closure, and
ontology-growth infrastructure landed in earlier passes
(\lpath{Logic/EmpiricalIntensionalFactorGraphBridge.lean},
\lpath{Logic/CredalConceptFixpointClosureBridge.lean},
\lpath{Logic/CredalConceptOntologyGrowthBridge.lean}).

\subsection{Lower-credal inheritance table}\label{sec:inheritance-table}

The credal inheritance table mirrors the precise version at line 1644 but operates
on \lid{LowerFormedConcept}:

\begin{definition}[\lid{lowerCredalConceptInheritanceTable},
\lpath{EmpiricalIntensionalFactorGraphBridge.lean:1778}]\label{def:lower-table}
\begin{lstlisting}
noncomputable def lowerCredalConceptInheritanceTable
    (Γ : Gate → EvidenceGate Q)
    (M : Obj → Attr → Q)
    (subConcept superConcept : LowerFormedConcept Γ M) :
    FiniteWitnessFeatureTable :=
  Interpretation.toFiniteWitnessFeatureTable
    (lowerFormedConceptInterpretation Γ M)
    subConcept superConcept
\end{lstlisting}
\end{definition}

\noindent Three exactness theorems follow: the prior at line 1787, the strength at
line 1804, and the log-ratio at line 1822. Each proves the credal interpretation's
read of the corresponding quantity equals the table read, via \lid{simpa} from the
underlying precise lemma. The packaged conjunction
\lid{lowerFormedConceptInheritance\_exact\_via\_table} at line 1985 is the credal
analogue of the line-1644 precise theorem.

\subsection{Closure transport}\label{sec:closure-transport}

The credal inheritance survives closure under WM-style rule sets:

\begin{theorem}[Threshold validity through closure,
\lpath{CredalConceptFixpointClosureBridge.lean:71}]
\label{thm:closure-transport}
Let $\mathit{seed}$ be a set of \lid{LowerFormedConceptPair}, $R$ a \lid{RuleSet}.
If a query-encoding $\mathit{encode}$ exposes the supported lower-formed
extensional-inheritance slice at strength level (the \lid{hEncode} hypothesis),
and every seed obligation is above threshold $\tau$ according to the exact generated
table, then the encoded seed query set is threshold-valid in the world model. The
companion \lid{leastRuleClosure\_thresholdValid\_of\_exactTableStrength} at line~114
shows the same obligations stay threshold-valid after applying any state-indexed
WM consequence rule set.
\end{theorem}

\subsection{Ontology-growth transport}\label{sec:growth-transport}

The credal inheritance is stable under ontology growth, provided the agreement
region is interaction-closed and the two MLN specs satisfy the
\lid{PaperUniformSmallTotalInfluence} (Dobrushin-style) bound:

\begin{theorem}[Stable under ontology growth,
\lpath{CredalConceptOntologyGrowthBridge.lean:87}]
\label{thm:growth-transport}
If a lower-formed-concept seed family is threshold-valid in MLN world-model state
$\mu_1$ and every encoded seed query is supported inside an agreement region
$\Gamma$, then the same seed family remains threshold-valid after ontology growth
to the agreeing spec on $\mu_2$. The companion line~145 theorem extends the result
through least-rule closure, and lines 212/280 push the result into perspectival
availability/admissibility.
\end{theorem}

\noindent Theorems~\ref{thm:closure-transport} and \ref{thm:growth-transport}
together demonstrate that the credal carrier is not a parallel stack. It integrates
with the existing world-model semantics: credal inheritance obligations transport
through both \emph{logical} closure (rule-set application) and \emph{ontological}
closure (Dobrushin-bound growth).
This is a conditional transport theorem: it does not assert that arbitrary
ontology growth is stable, but that stability follows under the stated
agreement-region and Dobrushin-style hypotheses.

\section{Credal truth coordinates and the Walley bridge}\label{sec:walley}

The credal layer reads $\STV{s}{c}$-like coordinates off the lower/upper envelopes
through definitions in \lpath{ProbabilityTheory/ImpreciseProbability/ProjectiveCredal.lean}.

\subsection{The three coordinates}

\begin{definition}[Credal envelope coordinates,
\lpath{ProjectiveCredal.lean}]\label{def:coords}
\begin{lstlisting}
noncomputable def credalEnvelopeWidth
    (C : CredalPrevisionSet Ω) (X : Gamble Ω) : ℝ :=
  upperEnvelope C X - lowerEnvelope C X

noncomputable def credalEnvelopeWidthComplement
    (C : CredalPrevisionSet Ω) (X : Gamble Ω) : ℝ :=
  1 - credalEnvelopeWidth C X

noncomputable def credalEnvelopeMidpoint
    (C : CredalPrevisionSet Ω) (X : Gamble Ω) : ℝ :=
  (lowerEnvelope C X + upperEnvelope C X) / 2
\end{lstlisting}
\end{definition}

\noindent The docstrings are explicit about the intended PLN reading:
\lid{credalEnvelopeWidth} is the imprecision coordinate that
``PLN confidence-like displays compress'';
\lid{credalEnvelopeMidpoint} is the point-estimate coordinate that
``PLN strength-like displays compress.''

The display map used in this paper is the canonical midpoint /
width-complement projection for this credal reading of PLN truth values:
\[
  \STV{s}{c}\ \longleftarrow\
  \bigl(\,\text{\lid{credalEnvelopeMidpoint}}\ C\ X,\
        \text{\lid{credalEnvelopeWidthComplement}}\ C\ X\,\bigr).
\]
Singleton-recovery: when \lid{credalSetDetermines}\ $C\ X$ holds, the lower and upper
envelopes coincide and the midpoint reads the precise value while the
width-complement reads $1$, recovering the precise PLN truth value. This is the
``three slots of freedom'' picture of~\cite{pln-credal-degrees-of-freedom} viewed
from the concept-formation side.

\noindent The same display lane is now packaged for credal inheritance pairs in
\lpath{EmpiricalIntensionalFactorGraphBridge.lean}.  The gamble
\lid{credalInheritanceGamble} uses the independent gate index
\lid{Gate × Gate}, so it matches the existing predicate
\lid{credalInheritanceJudgment}: the subconcept and superconcept each need upper
support, but not necessarily under the same admissible gate.  The global readout
theorems \lid{globalEnvelopeWidth\_credalInheritanceGamble\_eq},
\lid{globalEnvelopeWidthComplement\_credalInheritanceGamble\_eq}, and
\lid{globalEnvelopeMidpoint\_credalInheritanceGamble\_eq} prove the intended PLN
coordinates: \lid{credallyImpreciseInheritance} gives width $1$ and
width-complement $0$, while robust precision gives width-complement $1$.  The
review-facing bundle is \lid{CredalInheritanceTruthCoordinateCrown}.

\subsection{The Walley conjugacy cluster}\label{sec:walley-cluster}

A four-theorem cluster lands the conjugacy $\overline{P}^\ast = -\underline{P}_{\text{nat. ext.}}(-\cdot)$
at four scales of the projective credal spec: cylinder, finite-cylinder, global,
finite-global. We quote the global one verbatim; the cylinder one follows the same
pattern at a smaller scale.

\begin{theorem}[Global upper-envelope/natural-extension conjugacy,
\lpath{ProjectiveCredal.lean:6871}]\label{thm:walley-global}
\begin{lstlisting}
theorem globalUpperEnvelopePrevision_conjugate_eq_naturalExtension
    (S : ProjectiveLocalCredalSpec Window Global)
    (hNonempty : S.hasCompatibleCompletion)
    (hBdd : ∀ X : Gamble Global,
      BddAbove ((fun P : PrecisePrevision Global => P X) ''
        S.projectiveLimitCredalSet))
    (X : Gamble Global) :
    (S.globalUpperEnvelopePrevision hNonempty hBdd).conjugate X =
      S.globalNaturalExtension X
\end{lstlisting}
\end{theorem}

\noindent The proof is four lines: invoke
\lid{upperEnvelopePrevision\_conjugate\_eq\_lowerEnvelope} at the credal-set level
on \lid{S.projectiveLimitCredalSet}, with the spec's nonempty/compatibility
hypothesis as input. The three siblings live at lines 6268 (cylinder), 6523
(finite-cylinder), and 7065 (finite-global), with the same proof shape, scaled
hypotheses, and matching siblings going the other direction
(\lid{globalNaturalExtension\_conjugate\_eq\_upperEnvelope}).

\paragraph{What the cluster buys.} Each conjugacy theorem is a Walley-canonical
compatibility lemma: the spec's packaged upper-envelope prevision (a thing of type
\lid{UpperPrevision}) and the spec's packaged natural extension (a thing of type
\lid{LowerPrevision}) are conjugate. This is the load-bearing identity tying the
exact-natural-extension calculus to the credal-set envelope calculus on the
projective-credal infrastructure. The credal coordinates of
Definition~\ref{def:coords}, when read off
\lid{S.projectiveLimitCredalSet}, are the standard Walley coordinates of the
natural extension.

\section{The Mizar benchmark}\label{sec:mizar}

We exhibit the first concrete benchmark: an extractor for Mizar Mathematical
Library articles, a generated context, and a credal witness with computed
coordinates.

\subsection{The extractor}\label{sec:mizar-extractor}

\lpath{scripts/concept\_ontology/extract\_mizar\_article\_context.py} (330 lines)
walks a Mizar source file and emits a JSON observation context. Rows are
classified into four kinds (\lid{definition}, \lid{registration}, \lid{notation},
\lid{theorem}); attributes are extracted from the article's mode/attribute declarations
and predicate uses. Output schema:

\begin{lstlisting}
{
  "article":    string,
  "attributes": [ string ],
  "rows": [
    { "object":   string,
      "kind":     "definition"|"theorem"|"registration"|"notation",
      "index":    int,
      "preview":  string,
      "counts":   { attribute -> int } } ]
}
\end{lstlisting}

\noindent The extractor is deliberately small and conservative: it does not interpret
Mizar's full grammar, only enough to populate the attribute-count cells. The output
is then loaded by a generated Lean module
(\lpath{Generated/MizarConlat1.lean}) into the
\lid{BinaryFcaBenchmarkContext} surface from \lpath{BenchmarkControl.lean}.

\subsection{Working slice: \texorpdfstring{\lid{conlat\_1.miz}}{conlat\_1.miz}}\label{sec:mizar-data}

\lid{conlat\_1.miz} (Wojciechowski, 2002, on Concept Lattices, 3260 source lines)
is chosen for two reasons: it is a concept-lattice article, so the predicate
vocabulary is FCA-aligned, and it is well-scoped (mid-size, no
heavy preface). The extractor produces 14 rows over 9 attributes:

\begin{center}\small
\begin{tabular}{ll}
\toprule
\textbf{Attribute} & \textbf{Role} \\
\midrule
\lid{Attribute}            & FCA primitive \\
\lid{AttributeDerivation}  & Galois operator (attribute side) \\
\lid{ContextStr}           & 2-sorted context structure \\
\lid{FormalContext}        & Article's main concept-context object \\
\lid{Function}             & Underlying carrier function \\
\lid{Information}          & Incidence relation \\
\lid{ObjectDerivation}     & Galois operator (object side) \\
\lid{Subset}               & Stratification primitive \\
\lid{is-connected-with}    & Incidence predicate alias \\
\bottomrule
\end{tabular}
\end{center}

\noindent Row-kind histogram: 6~definitions, 5~theorems, 2~registrations, 1~notation.

\subsection{Toy control: the flying-family example}\label{sec:mizar-control}

\lpath{ConceptOntology/BenchmarkControl.lean} provides a small toy concept-context
(animals with traits) used to exercise the FCA recovery layer with concrete witnesses.
The positive witness
\lid{flyingFamilyConcept\_mem\_exact} at line 238 shows the flying-family concept is
in the exact concept family; the negative witness
\lid{batOnlyFlyingConcept\_not\_mem\_exact} at line 246 shows a non-closed pair is
correctly rejected; the lower-credal exclusion
\lid{flyingFamilyConcept\_not\_mem\_lower} at line 265 shows the lower-credal carrier
correctly rejects the flying-family concept under a strict-gate family that requires
two supporting tokens. This is the controlled benchmark surface against which the
Mizar layer is instantiated.

\subsection{The credal witness on \texorpdfstring{\lid{conlat\_1}}{conlat\_1}}\label{sec:mizar-witness}

\lpath{ConceptOntology/MizarBenchmark.lean} instantiates the
\lid{BinaryFcaBenchmarkContext} surface on the extracted \lid{conlat\_1} context.
A two-element gate family
\lid{MizarGate} := \{\lid{loose}, \lid{strict}\} is defined:
\lid{loose} = the exact (positive-support) gate; \lid{strict} = the
positive-threshold-2 gate (requires two supporting tokens per cell).

The article's \lid{ObjectDerivation} attribute concept is encoded as
\lid{objectDerivationLooseConcept} (line 67). The credal witnesses are:

\begin{itemize}
\item \lid{objectDerivationLooseConcept\_mem\_upper} (line 127): the concept is
in \lid{upperMizarConceptFamily} (it forms under the \lid{loose} gate).
\item \lid{objectDerivationLooseConcept\_not\_mem\_strict} (line 131): it is not
closed under the \lid{strict} gate's incidence relation.
\item \lid{objectDerivationLooseConcept\_not\_mem\_lower} (line 138): therefore it
is not in \lid{lowerMizarConceptFamily}.
\end{itemize}

\noindent The three computed coordinate theorems (lines 147, 160, 173) follow from
the upper-only-membership pattern:

\[
  \mathrm{width} = 1, \qquad
  \mathrm{midpoint} = \tfrac{1}{2}, \qquad
  \mathrm{width\text{-}complement} = 0.
\]

\noindent These are the maximum-imprecision coordinates one would expect for a
concept that lives in upper-but-not-lower under a binary gate family: the lower
envelope reads 0, the upper reads 1, the midpoint is $\tfrac{1}{2}$, and the
width-complement (i.e., PLN-style confidence) is at zero.

\subsection{Honest accounting}\label{sec:mizar-honest}

The reader should not over-interpret the Mizar witness. Specifically:

\begin{enumerate}
\item \textbf{One article, hand-picked gates.} The witness is computed on a single
\lid{.miz} file with two artificial gates (positive-threshold-1 versus
positive-threshold-2). The credal coordinates are not measured against natural
gate variation in the Mizar Library; they are demonstrated against codex-chosen
thresholds.
\item \textbf{$\mathrm{width} = 1$ is structural, not empirical.} Under a binary
gate family where one gate says yes and the other says no, the lower envelope reads
0 and the upper reads 1, giving width exactly 1. The witness verifies the calculus,
not the substantive credal claim about MML.
\item \textbf{External FCA checking is a separate gate.} The order-isomorphism
Theorem~\ref{thm:fca-order} is the theorem-side claim that our exact lane is
classical FCA. An empirical cross-check against ConExp or LatticeMiner on the
extracted context is straightforward (the JSON schema is convertible to standard
\lid{.cxt} format) but is not the load-bearing result of this paper.
\item \textbf{Gate-family choice is itself a research question.}
The curated-family pilot in
\lpath{artifacts/concept\_ontology/mizar\_family/credal\_threshold\_summary.json}
now runs threshold gates over 13 lattice/FCA-adjacent MML articles. In that pilot,
all 13 articles have more than one induced concept family and at least one
upper-minus-lower concept under the sampled count thresholds; the aggregate
unstable-concept count is 917. This is evidence that the extraction and credal
scoring harness is live, not evidence that the sampled threshold gates are the
right scientific gate family. On classical binary FCA contexts
(\lid{livingbeings\_en.cxt}, \lid{animals\_en.cxt}), the same threshold gates
degenerate: threshold $\geq 1$ recovers the original context and threshold $\geq 2$
yields a near-null one, leaving the credal envelope structureless. Row-drop and
column-drop \emph{observation} gates do produce meaningful lower/upper families on
those binary contexts. The diagnostic is that the credal layer applies not whenever
``we have a dataset'' but whenever ``we have a meaningful family of admissible
gates for that dataset.'' Choosing the gate family to match the evidence type
is part of the empirical work; the theory does not automate it.
\end{enumerate}

\noindent What this benchmark \emph{does} demonstrate is that the formal
infrastructure runs end-to-end on real Mizar source: extractor $\to$ context
$\to$ credal layer $\to$ coordinates, and that the same machinery can now be run
as a small article-family pilot. It is the first stone of a real benchmark,
not yet a benchmark result. \S\ref{sec:roadmap} discusses the experiments that
would turn this into one.

\subsection{The semantic-feature lane: status}\label{sec:mizar-semantic}

Roadmap item~(f) below --- extracting \emph{semantic} rather than lexical features
--- is now operational, and reporting its current state at true weight is more
useful than deferring it. For a curated slice of the \lid{LATTICE2} article
(\lid{MPTP}/\lid{chainy} export, theorems \lid{t14}--\lid{t25}) we extract, per
theorem, a Craig interpolant on the entailment $\Gamma \vdash \varphi$ (Vampire
symbol elimination) and a finite model of the satisfiable premise fragment (Vampire
FMB~\cite{vampire-fmb}); the predicate/function symbols of each artifact become the
attributes of an FCA context, and the credal gate family runs as in
\S\ref{sec:mizar-witness}. Three findings, each stated at its true weight.

\emph{(1) A reusable schema synthesizer.} A single recognizer-based procedure
replaces the earlier per-cluster extractors: given a theorem cluster it emits a
parameterized schema (a theorem-with-blanks) or refuses with ``no common schema.''
On \lid{t14}--\lid{t17} it recovers the left-identity/extremum schema --- \emph{in a
bounded lattice, the unique left identity of meet (resp.\ join) is the top (resp.\
bottom)} --- with the operation and the extremal element as the two parameter slots,
and it \emph{refuses} the mixed \lid{t14}--\lid{t19} cluster rather than fabricating
a common template. We stress that this schema is an elementary lattice fact: it
witnesses that the engine recovers a genuine unification from raw first-order export,
and is \emph{not} offered as a mathematical result. On the strip-mined tail
(\lid{t20}--\lid{t25}) the synthesized families drift toward library
\emph{encoding} regularities (self-relations of the \lid{BINOP\_1} operation
functor) rather than lattice mathematics --- a faithful signal that elementary
\lid{LATTICE2} is exhausted as a source of interesting structure.

\emph{(2) A cross-role bridge.} The cleanest non-degenerate concept is the pair
\lid{t15}+\lid{t18} and its dual \lid{t17}+\lid{t19}: a
left-identity-uniqueness theorem and a unity-identification theorem that target the
\emph{same} extremal object through the \emph{same} selector. This is a readable
concept-family bridge between two theorem forms, confirmed by direct FCA closure,
and is not a solver-timeout artifact.

\emph{(3) No gate yet perturbs the mathematics.} We ran three gate families.
Premise-neighborhood subsetting is a no-op (the prover already uses only the
neighborhood). Whole-symbol operator ablation is destructive: the resulting
instability is dominated by \lid{failure:timeout} --- proof \emph{failure}, not
changed content. A non-deletion ``right-partition'' family --- moving a support
record across the Craig interpolation cut instead of deleting it --- is
solver-robust (zero timeouts) and does change the interpolant vocabulary, but it
perturbs the \emph{feature} (which symbols are forced into the shared interpolation
vocabulary), \emph{not the theory}: every gate leaves the theorem's truth untouched.
It is a feature-stability diagnostic, not a conjecture generator. The gate family
that would generate conjectures is stated as item~(f).

\section{What real experiments would look like}\label{sec:roadmap}

We are explicit about which experiments would move the benchmark from
infrastructure-witness to scientific result.

\paragraph{(a) External FCA cross-check.} Convert
\lid{Generated/MizarConlat1.json} to a standard FCA \lid{.cxt} file, compute
the concept lattice in ConExp~1.1.4 or FCAbedrock, and compare against the
Lean-computed exact concept family for size, edges, and selected concept
witnesses. The order-isomorphism theorem predicts perfect agreement. ${\sim}2$ days
of engineering. This is sanity, not novelty.

\paragraph{(b) Multi-article scale.} A 13-article lattice/FCA-family pilot is now
materialized in \lpath{artifacts/concept\_ontology/mizar\_family/}: the extractor
writes per-article JSON and \lid{.cxt} files, the threshold-gate summary records
non-trivial upper/lower splits in every article of the pilot, and the aggregate
unstable-concept count is 917. The next scale target is $50$--$100$ MML articles.
The mathematical library has ${\sim}1500$ articles total; the
\lid{chainy} filter of MPTP~\cite{urban2003mptp} already produces well-scoped
subsets. The extractor's per-article extraction is fast; loading all articles into
one Lean module is not, so the harness needs an article-batched scoring layer rather
than a single monolithic context. The output is a real
$\mathrm{Concept} \times \mathrm{Article}$ presence matrix.

\paragraph{(c) Article-derived gates.} Replace the hand-picked threshold gates of
\S\ref{sec:mizar-witness} with article-derived gates (each MML article becomes one
admissible gate; the credal carriers then witness ``concepts robust across the
article portfolio'' versus ``concepts attested in at least one article''). This is
the move from synthetic to semantic credal variation.

\paragraph{(d) Galois-dual ghost hunt.} FCA theory predicts every closure operator
has a Galois dual; every concept-lattice node has a corresponding co-lattice node.
\lid{conlat\_1.miz} itself exhibits this: \lid{ObjectDerivation} and
\lid{AttributeDerivation} are Galois-dual operators. If the multi-article scan
identifies concepts in \(\mathrm{Upper}\setminus\mathrm{Lower}\) where the FCA-dual concept is
expected but not formally registered in some articles, those candidates are
predictions of ``formally missing concepts'' in MML that the credal layer surfaces
via duality reasoning. The first narrow lexical scan,
\lpath{artifacts/concept\_ontology/mizar\_family/duality\_ghosts.json}, reports
three one-sided \lid{Top}/\lid{Bottom} duality-mention candidates
(\lid{conlat\_1}, \lid{yellow\_2}, \lid{waybel\_0}); this is a triage list, not a
semantic confirmation of missing concepts. Verifying a small number of such
predictions against domain experts (Mizar maintainers, recent MML changelog) would
constitute a published discovery in the FCA-on-formal-math territory. We give this
a ${\sim}65\%$
probability of being feasible at the formalization layer; ${\sim}50\%$ of
producing at least one publishable discovery.

\paragraph{(f) Conjectures live in the prove--refute gap.} The semantic lane of
\S\ref{sec:mizar-semantic} is operational but has not yet surfaced a non-trivial
conjecture, because every gate built so far either leaves the theory fixed
(right-partition) or destroys provability without establishing falsity (operator
ablation). The gate that \emph{generates} conjectures is hypothesis weakening with a
dual oracle: drop a hypothesis from a synthesized schema, then ask both a prover
(does the weakened schema still prove?) and a model builder (does its negation admit
a finite countermodel?). Three outcomes partition the result: \emph{proved} (the
hypothesis was unnecessary --- a generalization theorem); \emph{refuted} (the
hypothesis was essential --- a sharpness witness); and \emph{neither, in budget} ---
the genuine open conjecture. The credal layer ranks candidates by gate-stability,
and a companion Markov-logic world-model formalization (which proves
$\mathrm{queryStrength} = \mathrm{queryProb}$, i.e.\ the PLN strength a world model
assigns to a query equals its Markov-network probability) attaches a
truth-likelihood to the third class. The honest discriminating questions then are
whether the credal coordinates beat a syntactic baseline (ENIGMA-style symbol
features, raw overlap) and whether the truth-likelihood is \emph{calibrated} against
hammer-provability on decidable held-out conjectures. Crucially, this experiment
requires a substrate richer than elementary lattice theory --- where the
prove--refute gap is empty within budget --- such as the \lid{WAYBEL}
continuous-lattice articles, or cross-article schema transfer in which a schema
proven in one domain becomes a conjecture in another.

\paragraph{(g) Annotator-disagreement complement.} Walley-canonical credal
benchmarks come from real annotator disagreement: CIFAR-10H~\cite{cifar10h} (51
labels per image), ChaosNLI~\cite{chaosnli} (100 annotators per item). With
appropriate object-attribute-gate reframing (classes as attributes, annotators as
gates, $\sim 200$--$500$ stratified items), these become tractable Lean
benchmarks. This is the headline credal-empirical complement to the
formal-mathematical Mizar benchmark of (b)--(d). Out of scope for the present paper.

\section{Related work}\label{sec:related}

\paragraph{Classical FCA.} Wille's~\cite{wille1982} original triadic-context-free
treatment; the canonical reference Ganter--Wille~\cite{gantnerwille1999}; the
survey volume~\cite{gantnerstumemwille2005}. Mathlib's
\lpath{Mathlib/Order/Concept.lean} provides a Lean~4 formalization of the
crisp concept-lattice carrier. The present work's recovery theorem
(Thm.~\ref{thm:fca-order}) bridges our exact concept layer into that mathlib
artifact; we do not duplicate it.

\paragraph{Generalized FCA.} Belohlavek's fuzzy FCA~\cite{belohlavek1999}
replaces the crisp incidence with degrees in $[0,1]$. Djouadi and
Prade~\cite{djouadi2009} introduce interval-valued and possibilistic contexts.
Murai--Mizumoto--Tahani~\cite{murai1989} treat the topic from a possibility-theoretic
angle. Triadic FCA~\cite{wille1995} adds a context dimension. None of these works
formalize a conservative reduction to crisp FCA in dependent type theory, nor
integrate with Walley's natural-extension calculus through a projective credal
spec, nor connect to a PLN-style truth-value calculus.

\paragraph{Walley framework.} The standard reference is~\cite{walley1991}; the
modern handbook is~\cite{augustin2014}. Natural extension's role as the canonical
coherent extension and as the conjugate of the upper envelope is treated there.

\paragraph{Lean formalizations of imprecise probability.} The
\lpath{Mettapedia/ProbabilityTheory/ImpreciseProbability/ProjectiveCredal.lean}
module is a multi-thousand-line Lean~4 formalization of credal sets, lower/upper
previsions, natural extension, and projective credal specs, on which the present
paper relies. To our knowledge it is the largest extant Lean~4 formalization in
this area.

\paragraph{PLN.} Goertzel--Iklé--Ari--Heljakka~\cite{goertzel2008} introduce the
two-number truth value; the companion characterization
paper~\cite{pln-credal-degrees-of-freedom} formalizes its three slots of degrees
of freedom over the same projective credal spec used here. Sibling formalization
papers \lid{xiPLN.tex} and \lid{nuPLN.tex} establish the broader PLN-on-Mettapedia
program.

\paragraph{Mizar.} Bancerek et al.~\cite{bancerek2018} give the canonical reference
for the Mizar Mathematical Library; Urban~\cite{urban2003mptp} introduces MPTP.
Premise-selection work over MML/MPTP includes Jakub\r{u}v--Urban
ENIGMA~\cite{jakubuv2019enigma} and the present project's sibling
PLN-kNN-NB~\cite{pln-knn-nb}. The Mizar extractor presented here is, to our
knowledge, the first FCA/credal-layer formalization of MML article structure;
prior MML-engagement work has been ATP-centric.

\paragraph{Automated theory formation and theory exploration.} The
concept-formation $\to$ conjecture $\to$ prove $\to$ counterexample cycle is not
new. Colton, Bundy, and Walsh's HR~\cite{colton1999hr} forms concepts and
conjectures over data tables and dispatches them to a prover and a model generator,
scored by deterministic interestingness measures; Colton and
Wagner~\cite{colton2007fca} applied FCA to mathematical discovery over the OEIS and
Graffiti conjectures. Theory-exploration systems QuickSpec~\cite{quickspec} and
Hipster~\cite{hipster} synthesize and filter equational conjectures for proof
assistants. Our \emph{cycle shape} and our \emph{use of FCA for mathematical
discovery} are therefore not the novel elements, and we do not claim them as such.
The delta is narrow and specific: a Walley-style credal lower/upper family
\emph{indexed by a gate family} over \emph{prover-derived} semantic artifacts
(Craig interpolants, finite models), with a Lean-4 exact/credal foundation and a PLN
truth-coordinate reading. We likewise claim no novelty for uncertainty-aware FCA in
general (cf.\ the generalized-FCA paragraph) nor for automated conjecturing in
general.

\section{Conclusion and open work}\label{sec:conclusion}

\paragraph{Recap of contributions.} (i) Type-theoretic credal carriers for
concept formation as proper subtypes; (ii) FCA recovery as a Lean~4
order-isomorphism theorem against mathlib's concept-lattice infrastructure;
(iii) singleton-collapse conservative-extension theorems for both lower and upper
carriers; (iv) integration with Walley's natural extension via a four-scale
conjugacy cluster on a projective credal spec under the compatible-completion
and boundedness hypotheses of \S\ref{sec:walley-cluster}; (v) a first-stone Mizar benchmark
with a real credal witness on \lid{conlat\_1.miz}'s \lid{ObjectDerivation} concept.

\paragraph{Open work.} The upper-side inheritance table and the basic two-sided
credal inheritance predicates have now landed:
\lid{upperCredalConceptInheritanceTable},
\lid{credalInheritanceJudgment}, \lid{credallyPreciseInheritance}, and
\lid{credallyImpreciseInheritance}. These predicates now have a PLN
truth-coordinate package through \lid{CredalInheritanceTruthCoordinateCrown},
including midpoint and width-complement readouts for inheritance. A 13-article
Mizar-family pilot and first narrow duality-ghost scan are now materialized as
artifacts. The remaining work is to scale the pilot to a larger, less curated MML
slice, replace convenience threshold gates with article-derived semantic gates, and
validate any ghost candidates against Mizar-domain expertise.

\paragraph{What the paper is not.} The paper is not a benchmark result. It is the
formal-infrastructure writeup of a layer that can support a benchmark result. The
load-bearing scientific question --- ``can the credal layer surface
formally-missing concepts in real mathematical libraries?'' --- requires the
multi-article scale of \S\ref{sec:roadmap}(b)--(d) and is left open.

\paragraph{Hygiene.} All theorems quoted in this paper are verbatim from the
public Lean~4 source as of the date listed on the title page. A grep across the
cited files (\lpath{ConceptOntology/CredalFormation.lean},
\lpath{ConceptOntology/FCARecovery.lean},
\lpath{ConceptOntology/BenchmarkControl.lean},
\lpath{ConceptOntology/MizarBenchmark.lean},
\lpath{EmpiricalIntensionalFactorGraphBridge.lean},
\lpath{CredalConceptFixpointClosureBridge.lean},
\lpath{CredalConceptOntologyGrowthBridge.lean},
\lpath{ProjectiveCredal.lean}) for \lid{sorry}, \lid{admit}, \lid{axiom}, and
\lid{by trivial} returns zero hits. Axiom checks on the load-bearing theorems
return the mathlib baseline
\{\lid{propext}, \lid{Classical.choice}, \lid{Quot.sound}\}.

\begin{thebibliography}{99}

\bibitem{wille1982}
R.~Wille.
\newblock Restructuring lattice theory: an approach based on hierarchies of concepts.
\newblock In \emph{Ordered Sets} (I.~Rival, ed.), Reidel, 1982, pp. 445--470.

\bibitem{gantnerwille1999}
B.~Ganter and R.~Wille.
\newblock \emph{Formal Concept Analysis: Mathematical Foundations}.
\newblock Springer, 1999.

\bibitem{gantnerstumemwille2005}
B.~Ganter, G.~Stumme, and R.~Wille (eds.).
\newblock \emph{Formal Concept Analysis: Foundations and Applications}.
\newblock LNAI 3626, Springer, 2005.

\bibitem{wille1995}
R.~Wille.
\newblock The basic theorem of triadic concept analysis.
\newblock \emph{Order}, 12(2):149--158, 1995.

\bibitem{belohlavek1999}
R.~Belohlavek.
\newblock Fuzzy Galois connections.
\newblock \emph{Mathematical Logic Quarterly}, 45(4):497--504, 1999.

\bibitem{djouadi2009}
Y.~Djouadi and H.~Prade.
\newblock Interval-valued fuzzy formal concept analysis.
\newblock In \emph{ISMIS 2009}, LNCS 5722, Springer, 2009, pp. 592--601.

\bibitem{murai1989}
T.~Murai, M.~Mizumoto, and N.~Tahani.
\newblock Possibility theory and concept lattices.
\newblock In \emph{Proc. IFSA 1989}, 1989.

\bibitem{walley1991}
P.~Walley.
\newblock \emph{Statistical Reasoning with Imprecise Probabilities}.
\newblock Chapman \& Hall, 1991.

\bibitem{augustin2014}
T.~Augustin, F.~P.~A. Coolen, G.~de Cooman, and M.~C.~M. Troffaes.
\newblock \emph{Introduction to Imprecise Probabilities}.
\newblock Wiley, 2014.

\bibitem{goertzel2008}
B.~Goertzel, M.~Iklé, I.~F. Goertzel, and A.~Heljakka.
\newblock \emph{Probabilistic Logic Networks: A Comprehensive Framework for
Uncertain Inference}.
\newblock Springer, 2008.

\bibitem{pln-credal-degrees-of-freedom}
Mettapedia Working Notes.
\newblock Credal degrees of freedom for PLN truth values.
\newblock Companion paper, this collection, 2026.

\bibitem{pln-knn-nb}
Mettapedia Working Notes.
\newblock PLN-kNN-NB on extended MPTP 5k.
\newblock Companion paper, this collection, 2026.

\bibitem{wmpln_concept_formation_guidance}
GPT-Pro 5.4.
\newblock WM-PLN concept formation grounding guidance.
\newblock \lpath{GPT-Pro\_Responses/WM-PLN/wmpln\_concept\_formation\_grounding\_guidance.md},
2026.

\bibitem{urban2003mptp}
J.~Urban.
\newblock MPTP -- Motivation, implementation, first experiments.
\newblock \emph{Journal of Automated Reasoning}, 33(3-4):319--339, 2004.

\bibitem{bancerek2018}
G.~Bancerek, C.~Bylinski, A.~Grabowski, A.~Kornilowicz, R.~Matuszewski, A.~Naumowicz,
K.~Pak, and J.~Urban.
\newblock The role of the Mizar Mathematical Library for interactive proof
development in Mizar.
\newblock \emph{Journal of Automated Reasoning}, 61(1-4):9--32, 2018.

\bibitem{jakubuv2019enigma}
J.~Jakub\r{u}v and J.~Urban.
\newblock Hammering Mizar by learning clause guidance.
\newblock In \emph{ITP 2019}, LIPIcs 141, 2019, pp. 34:1--34:8.

\bibitem{vampire-fmb}
G.~Reger, M.~Suda, and A.~Voronkov.
\newblock Finding finite models in multi-sorted first-order logic.
\newblock In \emph{SAT 2016}, LNCS 9710, Springer, 2016, pp. 323--341.

\bibitem{cifar10h}
J.~C. Peterson, R.~M. Battleday, T.~L. Griffiths, and O.~Russakovsky.
\newblock Human uncertainty makes classification more robust.
\newblock In \emph{ICCV 2019}, pp. 9617--9626.

\bibitem{chaosnli}
Y.~Nie, X.~Zhou, and M.~Bansal.
\newblock What can we learn from collective human opinions on natural language
inference data?
\newblock In \emph{EMNLP 2020}, pp. 9131--9143.

\bibitem{colton1999hr}
S.~Colton, A.~Bundy, and T.~Walsh.
\newblock Automatic concept formation in pure mathematics.
\newblock In \emph{IJCAI 1999}, pp. 786--791.

\bibitem{colton2007fca}
S.~Colton and D.~Wagner.
\newblock Using formal concept analysis in mathematical discovery.
\newblock In \emph{Calculemus/MKM 2007}, LNCS 4573, Springer, 2007, pp. 205--220.

\bibitem{quickspec}
K.~Claessen, N.~Smallbone, and J.~Hughes.
\newblock QuickSpec: guessing formal specifications using testing.
\newblock In \emph{TAP 2010}, LNCS 6143, Springer, 2010, pp. 6--21.

\bibitem{hipster}
M.~Johansson, D.~Ros\'en, N.~Smallbone, and K.~Claessen.
\newblock Hipster: integrating theory exploration in a proof assistant.
\newblock In \emph{CICM 2014}, LNCS 8543, Springer, 2014, pp. 108--122.

\end{thebibliography}

\end{document}
